% 3-implementation.tex starts here
\section{Implementation}\label{cap:implementare}

\subsection{Implementation overview}

The previous chapter established the technology choices that make up the \textbf{\gls{pan}} architecture (defined in \secref{subsec:pan-definitie}). In what follows I describe how this pattern translates into the actual structure of the repository: a single \gls{analog} application under \texttt{apps/}, a set of \gls{nx} libraries grouped by responsibility under \texttt{libs/}, and a unified execution chain through which a request travels from the \gls{analog} server-side runtime to the \gls{trpc} and \gls{prisma} layers.

% Keep the heading + the ~14cm file tree together: if less than 15cm is left on
% the current page, start them on the next one rather than splitting the tree.
\needspace{15cm}
\subsubsection{Structure of the Nx monorepo}\label{subsec:structura-monorepo}
\vspace{0.3cm}
\begin{TreeVerb}
kookio/
├── apps/
│   ├── web/                    # Analog application (Angular 21 + SSR),
│   │                             runs both in the browser and as
│   │                             a server runtime (Nitro)
│   └── web-e2e/                # Playwright suite for the critical
│                                 flows (login, onboarding, pantry,
│                                 planner)
├── libs/
│   ├── feature/                # user-facing flows
│   │   ├── auth/               #   login, registration, session
│   │   ├── onboarding/         #   profile setup steps
│   │   ├── pantry/             #   pantry management
│   │   ├── plan/               #   daily planner
│   │   ├── recipes/            #   recipe feed and details
│   │   ├── groceries/          #   derived shopping list
│   │   ├── saved/              #   saved collections
│   │   ├── dashboard/          #   main authenticated screen
│   │   └── shell/              #   layout, topbar, navigation
│   ├── prisma-client           # Prisma schema, migrations,
│   │                             generated Zod schemas, derived types
│   ├── server/                 # auth, logger, trpc
│   │                             (server runtime helpers)
│   └── shared/                 # ui, utility, mock
│                                 (components, helpers, fixtures)
├── nx.json                     # global Nx configuration
├── package.json                # main project manifest
├── tsconfig.base.json          # TypeScript config + path aliases
├── vercel.json                 # Vercel routes, headers and build
└── vitest.workspace.ts         # Vitest workspace (all projects)
\end{TreeVerb}
\vspace{0.3cm}

In practice, a change in \texttt{libs/feature/pantry} triggers only the checks for the affected library and for the modules that depend directly on it (via \texttt{nx affected}), not for the entire repository \cite{nx}. The topology between groups (\texttt{feature/} does not import from \texttt{feature/}, \texttt{shared/} depends on nothing internal) is enforced automatically by the \texttt{enforce-module-boundaries} rule, checked in \acrshort{ci}.

\subsubsection{End-to-end application flow}\label{subsec:flux-capcoada}

At a high level, a request goes through three layers: \textit{\Gls{analog} (\acrshort{ui} and \acrshort{ssr})} \(\rightarrow\) \textit{\acrshort{trpc} (typed API)} \(\rightarrow\) \textit{\gls{prisma} (data access)} \cite{analog,trpc,prisma_orm}. The transitions between layers do not go through an additional layer of manual serialization --- the \gls{typescript} types are shared from end to end, and input validation at the \gls{trpc} boundary uses \gls{zod} schemas derived from the \gls{prisma} schema \cite{trpc,zod,prisma_zod_generator}. The step-by-step breakdown of a concrete request (from the initial \acrshort{ssr} to the subsequent calls from the browser) can be found in \secref{subsec:flux-executie}, and the corresponding sequence diagram in \figref{fig:flux-executie}. The overall picture of the components and of the externally managed services is shown in \figref{fig:arhitectura}.

\begin{figure}[htbp]
  \centering
  \begin{tikzpicture}[
    font=\footnotesize,
    >={Stealth[length=2mm]},
    L/.style={draw, rounded corners, align=center, minimum height=8mm, inner sep=4pt},
  ]
    \node[L, fill=blue!5, minimum width=128mm] (client)
      {\textbf{Browser} \;---\; Angular 21 · Material · Tailwind · NgRx Signals};

    \node[L, fill=black!8, minimum width=30mm, below=20mm of client] (trpc) {tRPC · Zod};
    \node[L, fill=black!8, minimum width=34mm, left=8mm of trpc] (analog) {Analog / Nitro\\\acrshort{ssr} + serverless};
    \node[L, fill=black!8, minimum width=26mm, right=8mm of trpc] (prisma) {Prisma};
    \begin{scope}[on background layer]
      \node[draw, dashed, rounded corners, fit=(analog)(trpc)(prisma), inner sep=4mm] (vercel) {};
    \end{scope}
    \node[anchor=south west, font=\bfseries] at (vercel.north west) {Vercel};

    \node[L, fill=orange!18, minimum width=26mm, below=20mm of analog] (fb) {Firebase Auth};
    \node[L, fill=green!12, minimum width=22mm, below=20mm of trpc] (redis) {Redis};
    \node[L, fill=green!12, minimum width=28mm, below=20mm of prisma] (db) {CockroachDB};
    \node[L, fill=red!12, minimum width=20mm, right=8mm of db] (sentry) {Sentry};

    \draw[<->] (client.south) -- node[right=1pt]{HTTPS} (vercel.north);
    \draw[->] (analog.south) -- (fb.north);
    \draw[->] (trpc.south) -- (redis.north);
    \draw[->] (prisma.south) -- (db.north);
    \draw[->, dashed] (vercel.south east) -- (sentry.north);
  \end{tikzpicture}
  \caption[Overall architecture of the application]{Overall architecture:
    the browser (\gls{angular}) communicates over HTTPS with the \gls{vercel}
    platform, which hosts the \gls{analog}/\gls{nitro} runtime (\acrshort{ssr} +
    serverless functions), the \gls{trpc} layer (with \gls{zod} validation) and
    the \gls{prisma} client. Persistence and cross-cutting services are managed
    externally: \gls{cockroachdb} (data), \gls{redis} (cache,
    \textit{rate-limiting}), Firebase (authentication) and \gls{sentry} (error
    monitoring).}
  \label{fig:arhitectura}
\end{figure}

\subsection{Implementation of the application}

Whereas the previous section justified the technology choices, this one shows how each choice looks in the repository: which files materialize it, which conventions apply it and what concrete constraints it imposed at the code level. For each technology, the emphasis is no longer on \textit{why} it was chosen, but on \textit{what it actually produced} once integrated into \textit{\gls{kookio}}.

\subsubsection{Modern web architectures}\label{subsec:arhitecturi-implementare}

The \gls{randare} configuration of \textit{\gls{kookio}} is concentrated in \texttt{apps/web/vite.config.ts}: the \texttt{analog} plugin is called with \texttt{ssr: true}, with no prerendered routes. Thus, \acrshort{ssr} is the rendering mode for all the application's routes, including pages with stable content such as \texttt{/about}; \acrshort{ssg}, although supported by \gls{analog}, is not used, so that every response goes uniformly through the server runtime \cite{analog}.

On the \acrshort{ssr} path, the distinction between public routes (\texttt{/recipes}, \texttt{/legal/\allowbreak attributions}) and authenticated ones (\texttt{/pantry}, \texttt{/plan}) does not show up in the \gls{randare} strategy, but in what the server executes before serving the page: authenticated routes go through a \textit{guard} (route interceptor) that issues a redirect to \texttt{/login} if the Firebase session is missing, visible to \gls{playwright} as \texttt{ERR\_ABORTED} on Chromium and handled through \texttt{waitUntil: 'commit'} in the \acrshort{e2e} tests.

\subsubsection{Library topology and the absence of a repository layer}\label{subsec:topologie-biblioteci}

The application's \textit{feature} libraries are vertical slices: each one represents a complete user flow (\texttt{auth}, \texttt{onboarding}, \texttt{pantry}, \texttt{plan}, \texttt{recipes}) and contains the components, the state stores (NgRx Signals \textit{stores}) and the services that flow depends on. The topology enforced by the \gls{nx} \texttt{enforce-module-boundaries} rule is acyclic: \texttt{apps/} \(\rightarrow\) \texttt{libs/feature/} \(\rightarrow\) \texttt{libs/server/}, \texttt{libs/shared/}, \texttt{libs/prisma-client}, and \texttt{libs/feature/X} cannot import from \texttt{libs/feature/Y}. Communication between features happens through \texttt{providedIn: 'root'} state stores or through the \gls{trpc} procedures exposed by \texttt{libs/server/trpc} \cite{nx,nx_boundaries,trpc}. The project's dependency graph, collapsed onto the four \textit{scope} domains, is shown in \figref{fig:nx-graph}.

\begin{figure}[htbp]
  \centering
  \includegraphics[width=0.7\textwidth]{graph-all}
  \caption[The project's \gls{nx} dependency graph]{The project's dependency
    graph (\gls{nx}), collapsed onto the four \textit{scope} domains:
    \texttt{feature} (9 libraries), \texttt{server} (3), \texttt{shared} (4) and
    \texttt{prisma-client}, plus the \texttt{web} and \texttt{web-e2e}
    applications. The arrows show the dependency relationship; the topology is
    acyclic (\texttt{apps} \(\rightarrow\) \texttt{feature} \(\rightarrow\)
    \texttt{server}, \texttt{shared}, \texttt{prisma-client}), with the lower
    layers having no dependencies on the upper ones.}
  \label{fig:nx-graph}
\end{figure}

Organizing the feature libraries as vertical slices (each a complete user flow --- \texttt{auth}, \texttt{pantry}, \texttt{plan}), rather than as horizontal layers matching the processing stages, reflects Parnas's central conclusion: it is almost always incorrect to begin the decomposition of a system on the basis of a flowchart \cite{parnas1972}. Parnas contrasts the two modularizations of the KWIC system --- one by stages (input \(\rightarrow\) circular shift \(\rightarrow\) alphabetizing \(\rightarrow\) output), the other by information hiding --- and shows that in the first a change such as the storage format propagates into every module, whereas in the second it stays confined to a single one. Because design decisions transcend the time order of execution, modules do not correspond to processing steps. In \gls{pan}, this observation translates into the structural prohibition \texttt{feature/X} \(\not\rightarrow\) \texttt{feature/Y}: slices communicate through \gls{trpc} procedures or root-level state stores, not horizontally through stages.

The acyclic topology enforced by the \gls{nx} \texttt{enforce-module-boundaries} rule also corresponds to the \gls{uses} relation that Parnas requires to be a loop-free partial ordering \cite{parnas1972}. The benefits he attributes to this ordering can be found in \textit{\gls{kookio}}: the lower layers (\texttt{shared}, \texttt{server}, \texttt{prisma}) can be used without the upper ones, and “cutting off” the upper features leaves a reusable trunk --- exactly the property that allows a new feature library to be added incrementally without touching the rest of the application. It is important, however, to mark the distinction Parnas makes explicitly: a hierarchical structure (an acyclic \gls{uses} relation) and a “clean” decomposition (through information hiding) are two desirable but independent properties. The \gls{nx} \texttt{enforce-module-boundaries} rule guarantees the former --- the dependency graph stays acyclic ---, but does not by itself guarantee that each library actually hides a changeable design decision. The latter remains a discipline applied when drawing the boundaries: the \acrshort{ci} rule prevents the topology from degrading, not the absence of a real secret in a poorly delimited library.

The backend and frontend code coexist in the same repository: the \gls{prisma} client and the generated \gls{zod} schemas live in \texttt{libs/prisma-client}, and their types are imported directly into the \gls{angular} components \cite{prisma_zod_generator}. There is no classic object-oriented \textit{repository} layer between the \gls{trpc} procedures and \gls{prisma}; the routers call \texttt{ctx.prisma} directly, and the multi-step transactions --- including writes that touch several domains, such as the atomic transfer of bought items from the shopping list into the pantry (\texttt{groceries.addBoughtToPantry}, \secref{subsec:lista-cumparaturi}) --- use \texttt{prisma.\$transaction} \cite{trpc,prisma_transactions}. This direct coupling does not reduce the testability of the procedures: the \gls{trpc} context remains an injection point through which \texttt{ctx.prisma} can be replaced with a test double, a mechanism detailed in \secref{subsec:testare-validare}.

\subsubsection{Angular and server-side rendering}\label{subsec:angular-ssr-implementare}

The application's routes are defined based on the file structure (\textit{filesystem-based}) under \texttt{apps/\allowbreak web/\allowbreak src/\allowbreak app/\allowbreak pages/}: a file \texttt{about.page.ts} exposes \texttt{/about}, \texttt{cook.\allowbreak [recipeId].\allowbreak page.ts} exposes a dynamic route with a parameter, and a co-located \texttt{cook.\allowbreak [recipeId].\allowbreak server.ts} loads data before \gls{angular} renders the component \cite{analog,analog_routing}. Each \texttt{*.page.ts} is a \textit{default export} and is loaded on demand (\textit{lazy loading}) by the \gls{analog} router; there is no centralized routes file.

The application-level \acrshort{ssr} configuration is minimal: \texttt{app.config.ts} registers file-based routing, the HTTP client and hydration with replay of the events captured in the interval between \acrshort{ssr} and \gls{hydration} (the concrete mechanism is described in \secref{subsec:flux-executie}) \cite{analog}. Each import added to \texttt{app.config.ts}, however, carries a cost amplified by \acrshort{ssr}: the module is evaluated on the critical path of every route, and aggregating \textit{barrel} files (re-exports of the \texttt{export *} kind) are not eliminated by \gls{vite} in the development environment. The practical constraint is to import directly, through sub-paths, only the tokens actually used, leaving heavy components (Material datepicker, autocomplete) to be loaded from the level of the page that uses them.

\subsubsection{Data management with Prisma and CockroachDB}\label{subsec:gestionarea-datelor}

The database schema lives at \texttt{libs/prisma-client/schema.prisma}, alongside the migrations managed by \gls{prisma} and the \texttt{prisma-zod-generator} generator, which produces a validatable copy of each model in \gls{zod} form \cite{prisma_orm,prisma_zod_generator}. The complete schema --- the 20 models and the relationships between them --- is shown as an entity-relationship diagram in \anxref{anexa:model-date}. The derived types (initial data sets, form inputs, payloads for \gls{trpc} routes) are built by composition --- \texttt{.pick()}, \texttt{.omit()}, \texttt{.extend()} --- directly on top of the generated schemas, so that renaming a field in \texttt{schema.prisma} propagates automatically to the consumers \cite{prisma_zod_generator}.

Two structural decisions support this arrangement. First, the public \texttt{@kookio/prisma-client} entry point is safe to run in the browser: it exports only types and \gls{zod} schemas, while the \texttt{PrismaClient} instance is isolated under the \texttt{@kookio/prisma-client/server} sub-path, whose import is blocked by ESLint in frontend code \cite{prisma_orm}. Second, the \gls{prisma} instance is extended with an automatic \textit{retry-on-conflict} for the \gls{cockroachdb}-specific \textit{serializable} error codes (\texttt{P2034}, \texttt{P2036}, \texttt{40001} and the \textit{restart transaction} message); all write operations (\texttt{create}, \texttt{update}, \texttt{upsert}, \texttt{delete} and the \texttt{*Many} variants) go through this \textit{wrapper}, and the multi-step transactions are designed to be idempotent so that they can be retried automatically \cite{prisma_transactions,cockroachdb}.

The concrete connection to \gls{cockroachdb} is made through the \textit{driver adapter} pattern introduced in \gls{prisma} 7: the \texttt{PrismaClient} instance receives at construction an explicit adapter (\texttt{@prisma/adapter-pg}) that encapsulates the \texttt{pg} connection \textit{pool} for the \gls{postgresql} wire protocol \cite{prisma_orm,cockroachdb}. The connection pool is sized for the \textit{serverless} execution model: at most five connections per instance, with a connection \textit{timeout} of 10 seconds and a query \textit{timeout} (\texttt{statement\_timeout}) of 20 seconds, so that an ephemeral instance does not exhaust the database's connection quota and a stuck query does not hold a connection indefinitely. Starting with this version, \gls{prisma} drops the native engine written in Rust: the client runs entirely in JavaScript, and the database URL is moved from \texttt{schema.prisma} into a dedicated \texttt{prisma.config.ts} file. This separation between the typed client (generated from the schema) and the actual transport driver makes it possible, in principle, to change the data source without regenerating the client --- a flexibility that is useful for integration testing against a local database and for future migrations between \gls{postgresql}-compatible providers.

Compound uniqueness constraints are declared in \texttt{schema.prisma} and enforced at the database level --- for example \texttt{@@unique([userId, date, slotType])} on \texttt{MealSlot}, or the \texttt{userId}--\texttt{ingredientCatalogId} pair on \texttt{PantryItem} ---, and \gls{prisma} automatically generates the \textit{compound key} forms (of the kind \texttt{where: \{ userId\_ingredientCatalogId: \{ \ldots \} \}}) for the corresponding lookup queries \cite{prisma_orm}.

\subsubsection{Building the initial data set}\label{subsec:pipeline-date}

The data the database is populated with --- the catalog of ingredients, categories, culinary areas and the 594 recipes --- is not written by hand, but derived from a public source, \textit{TheMealDB}, through an offline fetching and cleaning pipeline. The source, however, is a \textit{crowd-sourced} database: the instructions are free text, the measures appear in heterogeneous forms (\texttt{"2 tbsp heaped"}, \texttt{"1/2 cup"}), and the ingredient names contain duplicates, plural forms and typos. Cleaning this data into a form that conforms to the \gls{prisma} schema requires deterministic, versioned and reproducible transformations --- which is why the pipeline is built as a set of \gls{typescript} scripts separate from the application, under \texttt{research/the-meal-db/}, each one consuming the artifact of the previous phase and producing a JSON \textit{snapshot}.

The pipeline runs in five phases. The \textit{fetching} phases query the source API --- the catalog (categories, areas, ingredients) and the recipes, the latter discovered by scanning the 26 letters of the alphabet, with up to 20 ingredient positions parsed for each recipe from the source's flat structure. The \textit{catalog cleaning} phase normalizes the names to a stable \texttt{snake\_case} key (diacritics removed, whitespace collapsed) used as a unique identifier throughout the pipeline, and auto-classifies each ingredient into one of the categories through an ordered list of twelve matching rules (\textit{first-match-wins}, with \texttt{PRODUCE} as the final catch-all rule; ingredients that do not match remain unclassified, for manual review). The \textit{recipe cleaning and enrichment} phase derives the structured fields (described below), and a final \textit{residual fixes} phase applies, directly to the canonical files, a set of idempotent mechanical rewrites.

One structural decision of the pipeline deserves highlighting: the separation between the automatic output and the canonical source. The cleaning scripts never overwrite the \texttt{recipes.json} file committed to the repository --- they write to an intermediate snapshot (\texttt{clean-recipes.json}), and promotion into the canonical source happens through a deliberate manual review step. As a result, hand-made refinements (re-segmentations of steps that the automatic parser missed, transcription fixes) survive re-running the pipeline, while purely mechanical fixes remain automated and idempotent. The source of truth is versioned, not regenerated on every run.

Three of the enrichment transformations are heuristics with an algorithmic role, not mere reformatting. Estimating the \textit{number of servings} maps each ingredient to its mass in grams --- volumes through the density of water, piece-type units through a table of per-piece weights --- excluding pure flavoring ingredients (salt, oil, vinegar) and negligible units; the total bulk mass is divided by a tunable constant of \(400\) g per serving --- an indicative weight for a main-course serving, chosen pragmatically and tunable, with no claim to nutritional value --- and clamped to the interval \([1, 12]\), with a default of \(2\) servings for recipes with fewer than two bulk ingredients or under \(200\) g of usable data. \textit{Difficulty} is a three-tier \textit{bucket} (\texttt{EASY}, \texttt{MEDIUM}, \texttt{HARD}) derived from a complexity score of \texttt{ingredients} \(+\ 2 \times\) \texttt{steps} --- steps weighted double, since they reflect effort, not just the shopping list --- relative to the median of the scores across the entire catalog: below \texttt{median} \(\times\ 0.8335\) is \texttt{EASY}, above \texttt{median} \(\times\ 1.1665\) is \texttt{HARD}, the rest \texttt{MEDIUM} (the bands are the median \(\pm\ 1/6\), so they yield an even split only for a uniform distribution; since recipe complexity is right-skewed, there are more \texttt{EASY} than \texttt{HARD} recipes). Finally, \textit{measure parsing} turns the free-text strings into \texttt{(quantity, \texttt{MeasureUnit}, notes)} triples --- ranges take the lower bound, and unrecognized units are moved into the notes field, without losing information.

The residual fixes phase is where the reality of crowd-sourced data shows most clearly: twenty-four distinct transformations, each idempotent (run a second time, it changes nothing). Rather than an enumeration, they fall into four classes of data quality issues: \textit{typographic and Unicode normalization} (typographic apostrophes and curly quotes brought to ASCII, Unicode fractions parsed, zero-width characters removed); \textit{referential integrity} (rewriting ingredient names from plural to singular according to the merge plan, so that recipes link correctly to the catalog, plus deduplicating ingredients repeated within the same recipe, which the \texttt{@@unique([recipeId, ingredientCatalogId])} constraint would reject); \textit{reconstruction of the step structure} (splitting steps that contained internally numbered sub-steps, removing header steps and metadata leaked as instructions); and \textit{regeneration of derived fields} (\textit{kebab-case} slugs from the title, with explicit mapping of extended Latin characters, and reclassifying difficulty after the step restructuring shifted the global median). Idempotence is guaranteed structurally --- each transformation matches only the uncorrected form or stops when the current state equals the target --- so the pipeline can be safely re-run after any manual refinement.

\subsubsection{Client–server communication with tRPC}\label{subsec:comunicare-trpc}

The \gls{trpc} routers are defined under \texttt{libs/server/trpc/src/lib/routers/}, one router per domain (\texttt{auth}, \texttt{pantry}, \texttt{onboarding}, \texttt{meal-plan}, \texttt{recipes}, \texttt{cooking-session} and others). Each procedure is composed from one of three \textit{builders}: \texttt{publicProcedure} (an alias for \texttt{procedure}, exposed anonymously, without authentication, defined in \texttt{libs/server/trpc/src/lib/trpc.ts}), \texttt{protectedProcedure} (requires a Firebase session verified in \texttt{createContext}) and \texttt{onboardedProcedure} (additionally requires \texttt{onboardingDone === true} on the user, otherwise it rejects the request with \texttt{PRECONDITION\_FAILED}), the last two declared in \texttt{libs/\allowbreak server/\allowbreak trpc/\allowbreak src/\allowbreak lib/\allowbreak middlewares/\allowbreak auth.ts} \cite{trpc,trpc_middleware}. The convention is a structural one: a \texttt{grep} for the builder used in a router immediately shows what surface the endpoint exposes, without having to inspect each locally composed \textit{middleware} layer.

On the client side, the \gls{trpc} instance is created in \texttt{apps/web/src/trpc-client.ts}: the Firebase \textit{ID token} is attached as the HTTP \texttt{Authorization} header through the \texttt{headers} option of the transport \textit{link} (a reactive \textit{signal} read on every request), and a dedicated \texttt{authAwareFetch} (i) proactively refreshes the token when it expires in less than five minutes, (ii) retries the request once with a new token on a \texttt{401} response, (iii) resolves relative URLs to absolute form in the \acrshort{ssr} context (Node.js has no implicit base for \texttt{fetch('/api/\ldots')}). In this configuration the \gls{trpc} procedures return RxJS \texttt{Observable}s, and the \gls{angular} consumers choose between \texttt{from(\ldots).pipe(\ldots)} for a reactive stream and \texttt{await firstValueFrom(from(\ldots))} for a single result \cite{trpc}.

\subsubsection{Auxiliary h3 routes and the Redis infrastructure}\label{subsec:rute-h3-redis}

Besides the \gls{trpc} handler (\texttt{apps/\allowbreak web/\allowbreak src/\allowbreak server/\allowbreak routes/\allowbreak api/\allowbreak trpc/\allowbreak [trpc].ts}), the application exposes two additional \texttt{h3} routes: \texttt{/api/health} for the \textit{healthcheck} and \texttt{/api/internal/bench-db} for the \gls{cockroachdb} latency benchmark (used in the evaluation section). A single \texttt{h3} \textit{middleware} (\texttt{apps/web/src/server/\allowbreak middleware/security.ts}) covers two cross-cutting requirements: an explicit \texttt{404} for static files that are not found (\texttt{/assets/*}, \texttt{/fonts/*} and others) --- to avoid the situation in which \acrshort{ssr} would return \texttt{200} with HTML for a \textit{bundle} replaced between deploys --- and rejecting with \texttt{404} the usual automated scanning patterns (\texttt{/wp-admin}, \texttt{.env}, \texttt{.git}, \texttt{.php} extensions). Security headers are applied at two levels. At the network \textit{edge}, through \texttt{vercel.json}, \gls{hsts} is applied with a one-year duration and \texttt{includeSubDomains}, which forces HTTPS access across the entire domain. The \gls{csp} policy is issued per request by a dedicated \gls{nitro} plugin (\texttt{apps/\allowbreak web/\allowbreak src/\allowbreak server/\allowbreak plugins/\allowbreak csp.ts}): for each HTML response, the plugin generates a unique \textit{nonce}, attaches it to the application's inline scripts and includes it in the header, keeping \texttt{script-src} strict (the own origin plus the nonce) without allowing the execution of arbitrary inline code; the remaining sources stay limited to the own origin, with a \textit{frame} exception for the Firebase authentication domain.

The Redis component lives in \texttt{apps/web/src/server/infra/} and is designed to tolerate unavailability. \texttt{getRedis()} is a singleton with a 60-second \textit{circuit breaker}: after a failed connection, the following calls \textit{fail fast} instead of recreating the client \cite{redis}. On the read path (\texttt{groceries.list}), a utility function silently catches Redis failures and returns the empty set, so that the page renders without the “bought” state; on the write path, the response includes a \texttt{persisted: false} field if Redis is unavailable, allowing the client to keep the \textit{optimistic update} and display a notification. Rate limiting (\textit{rate-limiting}, \texttt{rateLimitFixedWindow}) uses an \texttt{INCR + EXPIRE} fixed window and is applied both to the \textit{benchmark} route (10 requests per window) and to the \gls{trpc} handler: 10 requests for registration (\texttt{auth.register}) and 300 for the other procedures. In line with the unavailability-tolerance principle described above, rate limiting is \textit{fail-open} --- if Redis does not respond, the request is let through instead of being rejected, so that a failure of the \textit{cache} layer does not block authentication or writes. Also to protect the free-tier Redis instance, a weekly scheduled task (\textit{cron}) calls the \texttt{/api/cron/\allowbreak redis-keepalive} route, keeping it active to prevent the auto-suspension cycle.

\subsubsection{Logging and observability}\label{subsec:jurnalizare}

The \gls{logging} layer is isolated in \texttt{libs/server/logger/}: a single \gls{pino} instance is created with an environment-dependent transport --- \texttt{pino-pretty} with \textit{colorize} and a short timestamp in development, raw JSON in production (\gls{vercel}) \cite{pino}. On the production path, the instance applies an explicit list of redacted fields (\textit{redact list}: \texttt{password}, \texttt{*.token}, \texttt{authorization}) to prevent credentials from leaking into the logs.

Besides the global logging instance, \texttt{createLogger(namespace, bindings)} returns a child \textit{logger} with a \textit{namespace} label from a fixed set (\texttt{audit}, \texttt{http}, \texttt{trpc}, \texttt{prisma}) plus additional \textit{bindings} --- for example \texttt{createLogger('trpc', \{ router: 'pantry' \})} produces \texttt{ns: 'trpc'} alongside \texttt{router: 'pantry'} on every emitted entry ---, and \texttt{createAudit(bindings)} returns a logger specialized for audit events (with the \texttt{audit: true} flag applied automatically). Messages are always structured --- \texttt{logger.info(\{ userId, action \}, 'message')} --- not string interpolations, so that the resulting logs can be filtered directly on the aggregation platform.

On top of the logging layer, error monitoring in production is provided by \gls{sentry}, integrated symmetrically on the client and the server and inactive (\textit{no-op}) until an ingestion key (\acrshort{dsn}) is configured, so that development environments and deployments prior to the configuration remain untouched. On the client, a global \texttt{ErrorHandler} captures unhandled errors, and a \texttt{reportError} \textit{helper} isolated in \texttt{libs/shared/observability/} is called at the error-handling points in the \textit{stores} --- the \texttt{catchError} blocks that consume the error to turn it into interface state ---, making otherwise silent failures visible as well. On the server, the \texttt{wrap()} function in \texttt{libs/server/logger/} becomes the single capture point: the \texttt{error} and \texttt{fatal} methods forward the \texttt{err} field to \texttt{Sentry.captureException} before delegating to \gls{pino}, so that every \texttt{log.error(\{ err \})} call in the \gls{trpc} routers reports automatically, without instrumenting each error point. The \texttt{@sentry/node} client is initialized only once, at server startup, through a \gls{nitro} \textit{plugin} that also registers a \textit{hook} for unhandled errors, filtered to \texttt{5xx} codes so as not to report \texttt{4xx} client errors. The symmetry is deliberate: unhandled errors are caught by the global \textit{handlers} on both ends, and handled ones by the two single capture points --- \texttt{reportError} on the client, \texttt{wrap()} on the server. Besides errors, the client also collects performance transactions through a \textit{tracing} integration with a sampling rate of 10\,\%, enough to observe the pages' vital metrics without exceeding the free-tier quota. At build time, the \texttt{@sentry/\allowbreak vite-plugin} \textit{plugin} uploads the \textit{source maps} to \gls{sentry} and removes them from the delivered \textit{bundle}: generated in \textit{hidden} mode, without the \texttt{sourceMappingURL} comment, they are never served to the client, and each deployment is tagged with a \textit{release} equal to the Git \textit{commit} \textit{hash}, so that a production error points exactly to the \textit{deploy} that introduced it \cite{sentry}.

\subsubsection{Execution flow of a request}\label{subsec:flux-executie}

The flow summarized in \secref{subsec:flux-capcoada} unfolds in three stages.

\textbf{Stage~1 (initial rendering).} The HTTP request reaches the \gls{analog} server-side runtime (based on \gls{nitro}), and the \gls{angular} instance is executed on the server to produce the initial HTML \cite{analog,nitro}.

\textbf{Stage~2 (querying the persistence layer).} The \gls{angular} components trigger calls to the \gls{trpc} procedures of the server-side layer; each procedure validates its inputs with a \gls{zod} schema derived from the \gls{prisma} schema and uses \texttt{ctx.prisma} to query \gls{cockroachdb} \cite{trpc,zod,prisma_zod_generator,prisma_cockroachdb}. The results propagate back through the \gls{trpc} layer to the components that take part in the initial rendering.

\textbf{Stage~3 (hydration and subsequent interaction).} The rendered HTML reaches the browser, and the \gls{angular} instance takes over through \gls{hydration} --- the DOM tree produced on the server is reused, and the events captured in the meantime are replayed by \texttt{withEventReplay()} \cite{analog}. From then on, subsequent interactions call the \gls{trpc} procedures directly from the browser, without a new \gls{randare} on the server \cite{trpc}. \figref{fig:flux-executie} shows this flow as a sequence diagram.

\begin{figure}[htbp]
  \centering
  \resizebox{\textwidth}{!}{%
  \begin{sequencediagram}
    \newthread{br}{Browser}
    \newinst[2.5]{sv}{Analog / Nitro}
    \newinst[2.5]{tr}{tRPC}
    \newinst[2]{pr}{Prisma}
    \newinst[2]{db}{CockroachDB}

    \begin{call}{br}{HTTP GET /route}{sv}{initial HTML}
      \begin{call}{sv}{procedure (Zod validation)}{tr}{result (SuperJSON)}
        \begin{call}{tr}{ctx.prisma.<op>}{pr}{data}
          \begin{call}{pr}{SQL}{db}{rows}
          \end{call}
        \end{call}
      \end{call}
    \end{call}
    \postlevel
    \begin{call}{br}{direct tRPC call (post-hydration)}{tr}{response}
      \begin{call}{tr}{ctx.prisma.<op>}{pr}{data}
      \end{call}
    \end{call}
  \end{sequencediagram}}
  \caption[Sequence diagram of a request]{Sequence diagram of a
    request. \textbf{Initial rendering (\acrshort{ssr}):} the HTTP request reaches
    the \gls{analog}/\gls{nitro} runtime, which queries the persistence layer
    through the \gls{trpc} (with \gls{zod} validation) \(\rightarrow\) \gls{prisma}
    \(\rightarrow\) \gls{cockroachdb} layers and returns the initial HTML.
    \textbf{Subsequent interaction:} after hydration, calls go directly
    from the browser to \gls{trpc}, without a new server-side rendering.}
  \label{fig:flux-executie}
\end{figure}

\subsubsection{Testing and validation}\label{subsec:testare-validare}

The \gls{vitest} suite is orchestrated by \texttt{vitest.workspace.ts} at the root of the repository, which automatically collects all the \texttt{vite.config.*} and \texttt{vitest.config.*} files from the \gls{nx} projects \cite{vitest,vite}. The specification files (\textit{specs}, \texttt{*.spec.ts}) are co-located with the tested code, both in the \textit{feature} libraries and in the server libraries (for example \texttt{pantry.router.spec.ts} next to \texttt{pantry.router.ts}). Granularity is the guiding principle: \texttt{nx affected -t test} runs only the specs touched by the current changes.

At the level of the server libraries, the \gls{trpc} procedures are tested in isolation, without starting an HTTP server and without a real database: each spec builds a fake context in which \texttt{ctx.prisma} is a test double with \texttt{vi.fn()} methods declared exactly for the operations touched by the procedure, and the procedure is invoked through \texttt{appRouter.createCaller(ctx)} \cite{trpc,vitest}. Since the \gls{prisma} client reaches the procedures through the context, not through a direct import, the context takes over the isolation role that a \textit{repository} layer would have had: its absence (\secref{subsec:topologie-biblioteci}) does not sacrifice testability, it only moves the injection point into the context.

The \gls{playwright} suite lives at \texttt{apps/web-e2e/} and has its own \texttt{global-setup.ts} that warms up all the application's routes before the first \textit{worker} process starts --- a measure needed to avoid the time constraints at the first start of \gls{vite} in the development environment \cite{playwright}. The critical flows (authentication, registration, completing \textit{onboarding}, adding to the pantry, viewing the planner) go through the interface in Chromium; protected routes use \texttt{waitUntil: 'commit'} to avoid the \texttt{ERR\_ABORTED} error produced by \gls{angular} redirects during \gls{hydration}.

The discipline most visible in practice applies to refactorings that change the dependencies injected into a shared component: the tests of the library in question keep passing (the local \textit{stubs}, test stand-ins, are updated), but the specs of the dependent modules may fail with \texttt{NullInjectorError}. To detect these regressions, \texttt{nx affected -t test} is run, which follows the \gls{nx} dependency graph and runs everything that is affected \cite{nx}.

\subsubsection{Development and deployment}\label{subsec:dezvoltare-deploy}

Environment configuration is centralized in the \gls{vercel} dashboard, separately for \texttt{production}, \texttt{preview} and \texttt{development} \cite{vercel_cli}. Locally, \texttt{vercel env pull} synchronizes the values of the current environment into a \texttt{.env.local} (excluded from version control through \texttt{.gitignore}), removing the need for any secret to be versioned in Git.

Promotion between environments follows the \texttt{dev} \(\rightarrow\) \texttt{test} \(\rightarrow\) \texttt{main} flow. The \texttt{dev} branch feeds the development \textit{Preview} environment; the \texttt{test} branch triggers the \acrshort{e2e} suite on the testing \textit{Preview} environment before promotion to \texttt{main}, and the \texttt{main} branch feeds the production \textit{deploy} in the \texttt{fra1} region (Frankfurt) \cite{vercel_cli}.

The production build runs \texttt{nx build web} with \texttt{NODE\_OPTIONS=\allowbreak --max-old-space-size=6144}.
The default V8 heap limit is insufficient for bundling the entire \gls{angular} Material stack in the \gls{nitro} context: the phases that transform the abstract syntax tree (AST) and generate the \textit{source maps} accumulate large volumes of temporary objects, and the build fails with \texttt{JavaScript heap out of memory} before completing.
Raising the limit to 6\,GB provides the necessary headroom for both phases.
\gls{prismaclient} is regenerated in the automated \textit{pipeline} --- the \texttt{generate} target is a dependency of the \texttt{build} target ---, so that the \gls{typescript} client and the \gls{zod} schemas derived from \texttt{schema.prisma} are always in sync with the data model before the consumers are compiled \cite{prisma_orm}. Since \gls{prisma} 7 no longer ships a platform-dependent native engine (the client runs entirely in JavaScript through a \textit{driver adapter}), regeneration concerns only the typed code, not binaries specific to the execution environment.

\subsubsection{Architectural patterns discovered during implementation}\label{subsec:pan-patterns}

Besides the structural principles established \textit{a priori} (\secref{subsec:pan-definitie}), implementing the \textit{\gls{kookio}} application exposed three architectural patterns that were not predictable from the initial choices, but that consolidated into generalizable lessons for any \gls{pan} project. Unlike the intrinsic limitations of the stack (\secref{subsec:pan-limitari}), the patterns described here are reusable \textit{solutions}, codified in the project's guidelines repository to prevent the same problem from recurring:

\begin{itemize}
  \item \textbf{Authenticated SSR loaders require client-side hydration of authenticated data.} When first written, the \texttt{*.server.ts} loaders co-located with the dynamic routes called authenticated \acrshort{trpc} procedures directly through \texttt{appRouter.createCaller(ctx)}, assuming that the Firebase session from \texttt{createContext} is available both on the first \acrshort{ssr} and on subsequent client-side navigation. The assumption was false: on client-side navigation, \gls{analog} executes the \textit{loader} as an ordinary \texttt{fetch} that does not propagate the \texttt{Authorization} header with the Firebase \textit{ID token}, and the authenticated procedures respond with \texttt{401}. The correct pattern, codified in the repository: the \textit{loader} loads only public data (\textit{recipes.byId}, \textit{about page content}), while authenticated data (the availability of ingredients in the pantry, the “saved” state of a recipe, the user's preferences) is hydrated afterwards, from the browser, through the \gls{trpc} client with \texttt{authAwareFetch}. In terms of general applicability, the pattern comes down to the observation that \acrshort{ssr} and \textit{Bearer token} authentication do not coexist for free in any hybrid meta-framework --- any reader evaluating \gls{pan} for their own project should audit this dimension before relying on \acrshort{ssr} \textit{loaders} for sensitive data.

  \item \textbf{Route names in file-based routers must be semantic, not numeric.} When building the page for the \texttt{500} error code, the natural choice was a \texttt{/500} route (parallel to the already working \texttt{/404}). At runtime, the route did not activate: the server responded with \texttt{200} and the HTML of the home page. The cause, uncovered through gradual investigation, was a quirk of the \gls{analog} file-based router: \textit{all-numeric} segments are interpreted as invalid dynamic parameters, and the \textit{matcher} silently falls through the list of routes without signaling a configuration error. The correct pattern, codified in the repository: \textit{semantically} named routes (\texttt{/server-error}, \texttt{/not-found}, \texttt{/maintenance}) instead of numeric codes. In terms of general applicability, the pattern comes down to the observation that any file-based routing convention must separate \textit{numeric codes} (HTTP protocol artifacts) from \textit{route names} (semantic identifiers for the user) --- even if, in the world of protocols, the two seem interchangeable.

  \item \textbf{A single defensible heuristic for estimating quantities, two consumers.} The offline data cleaning pipeline (the \textit{seed} stage) and the runtime logic (pantry decrementing, shopping list derivation) initially consumed almost identical conversion heuristics, written separately --- a manual \texttt{INGREDIENT\_PIECE\_GRAMS} table at \textit{seed} time and their own \textit{carve-outs} at runtime ---, which forked every extension and, over time, produced behavioral divergences. Consolidating the two into a single source of truth (the pair of pure functions \texttt{toGrams}/\texttt{fromGrams} and the \texttt{IngredientCatalog.pieceGrams} column in the database schema) is detailed in \secref{subsec:algoritm-grame}. In terms of general applicability, the pattern comes down to a design rule: \textit{one defensible heuristic, two or more consumers, never one heuristic per consumer}. A critique of the estimate (for example, “the density of water is a poor approximation for oils”) applies to all consumers at once; an improvement propagates the same way.
\end{itemize}

The three patterns share a common structure: each emerged from \textit{implementation friction} (the \acrshort{ssr} \textit{loader} that responded with \texttt{401}, the \texttt{/500} route that did not activate, two almost identical heuristics written separately), and the solution was not adding functionality, but \textit{repositioning} the already existing piece at a level closer to the real invariant of the problem. From the standpoint of architectural auditing, this discovery pattern --- \textit{friction} \(\rightarrow\) \textit{lesson} \(\rightarrow\) \textit{codification} --- has been observed in the software engineering literature as a characteristic of evolutionary systems, in which the correct invariant is learned from implementation experience, not imposed by the initial design.

\subsection{Functional flows and product algorithms}\label{subsec:fluxuri-functionale}

The previous sections described how the architectural layer --- \gls{analog}, \acrshort{trpc}, \gls{prisma} --- materializes the \gls{pan} architecture at the code level. This section goes one level lower, to the \textit{product algorithms} that realize the application's value proposition: progressive filtering of recipes based on the pantry inventory, transactional decrementing of quantities after cooking, and automatic derivation of the shopping list from the planned slots. The three flows converge into a single \textit{estimation algorithm} that unifies the data cleaning pipeline (described in \secref{subsec:gestionarea-datelor}) with the production runtime logic.
Progressive filtering of recipes based on the pantry inventory (\secref{subsec:filtru-progresiv}) and automatic derivation of the shopping list (\secref{subsec:lista-cumparaturi}) operationalize two of the prevention strategies identified by Schanes et al.~\cite{schanes2018}: cooking with what is already stored at home and checking the inventory before shopping.

\subsubsection{Progressive recipe filtering}\label{subsec:filtru-progresiv}

On the \texttt{/recipes/match} route, the user sees the recipe catalog partitioned into three levels, labeled in the interface \textit{Alright} (all non-negligible ingredients are present in the pantry), \textit{Almost} (one to three non-negligible ingredients are missing) and \textit{All} (all public recipes, unfiltered). The \gls{trpc} procedure \texttt{recipes.match} receives a \texttt{level} field with exactly these three values (\texttt{'Alright'}, \texttt{'Almost'}, \texttt{'All'}), along with an optional \texttt{cursor} and a \texttt{limit} (default 20, maximum 50), and returns a page of recipes with the matching metadata the interface needs: the number of missing ingredients (\texttt{missingCount}), the names of the first three of them (\texttt{missingNames}) and the recipe's total number of ingredients (\texttt{totalIngredients}). The page also carries a \texttt{nextCursor} for pagination and a \texttt{partialScan} flag used by the client to automatically continue a scan interrupted by the examination budget.

The algorithm is, in essence, a \textit{set difference} between the recipe's ingredients and the contents of the user's pantry. For each recipe in the \texttt{PUBLIC} catalog, the procedure iterates through the list of \texttt{RecipeIngredient}s and skips the ingredients with a negligible unit (\texttt{TO\_TASTE}, \texttt{TO\_SERVE}, \texttt{GARNISH}, \texttt{PINCH}) --- salt, pepper, a parsley garnish and oil cannot block a recipe from the \textit{Alright} level just because they are missing from the explicit pantry records. The remaining ingredients are looked up in a map built from \texttt{pantryItem.findMany(\{ where: \{ userId \} \})}, compared at the level of \textit{existence} (the presence of the ingredient in the pantry, regardless of quantity). In the version delivered with this thesis, the quantity comparison is not performed; a rare false positive is preferred over unit logic that would duplicate the pipeline described in \secref{subsec:algoritm-grame}.

The decision to keep the matching logic simple --- \textit{presence}, not \textit{sufficient quantity} --- was taken deliberately: the cost of a false positive is low (a recipe appears as cookable although the user has an insufficient quantity of an ingredient), whereas the quantitative variant would duplicate the unit conversion logic of the grams pipeline (\secref{subsec:algoritm-grame}) at the matching boundary, with no proportional gain in perceived accuracy. For future development, the quantitative variant is built directly on top of that pipeline --- ingredients are normalized to grams before comparison, and the \textit{Alright} level requires the pantry to cover the recipe's consumption in grams for all non-negligible ingredients.

\subsubsection{The cooking session and transactional decrementing}\label{subsec:sesiune-gatit}

On the \texttt{/cook/:recipeId} route, the user goes through the recipe's steps and, at the end, a dialog offers a list of pre-filled quantities (proportional to the selected number of servings) for confirmation. On submission, the client calls \texttt{cookingSession.complete}, which decrements the pantry in a single \gls{prisma} transaction (\texttt{\$transaction}). For each confirmed \texttt{(ingredientCatalogId, consumedQuantity)}, the procedure looks up the corresponding row by the compound key \texttt{(userId, ingredientCatalogId)} (the \texttt{@@unique} constraint declared in \secref{subsec:gestionarea-datelor}), subtracts the consumed quantity from the existing quantity and --- this is where the estimation algorithm described in the next section comes in --- deletes the row if the remaining quantity drops below a minimum threshold (\texttt{0.5} g, the noise threshold of a pinch of salt). The entire operation runs atomically: either all the decrements go through and \texttt{cookingSession.completedAt} is updated, or nothing is applied.

Two implementation details are worth mentioning. First, the decrementing is done internally in grams (not in the display unit of the pantry row), to allow consumption in units different from the storage unit --- for example, a recipe calls for \texttt{200 g} of flour, but the user entered the flour as \texttt{0.5 kg}. The conversion back to the display unit happens at the end, before \texttt{prisma.update}. Second, if the user indicated a consumption greater than the quantity available in the pantry (an input mistake, or a recipe with a larger quantity entered later), the quantity is clamped to zero and the row is deleted, instead of receiving a negative value. These robustness behaviors at the edge of the data flow were validated through eighteen unit specification cases for \texttt{cookingSession.complete} at the time of the grams pipeline refactoring (28 May 2026).

\subsubsection{The grams pipeline: one heuristic, two consumers}\label{subsec:algoritm-grame}

The two previous flows --- progressive filtering (\secref{subsec:filtru-progresiv}) and the cooking session (\secref{subsec:sesiune-gatit}) --- depend on the same practical question: \textit{how do we convert quantities between heterogeneous units?} The recipe calls for \texttt{1 cup} of sugar and the pantry holds \texttt{500 g}; the recipe calls for \texttt{2 lemons}, but the pantry has only \texttt{1 piece}. A classic approach would separate each case (\textit{carve-out}): a converter for weights, a converter for volumes, a special rule for pieces. Over time, such layering produces branches that grow with every unit added.

The solution adopted in \textit{\gls{kookio}} is a pair of pure functions --- \texttt{toGrams} and \texttt{fromGrams} --- that normalize all quantities to grams before any operation and convert back to the display unit only at the end. The pipeline handles four classes of units --- weight, volume, pieces and taste/garnish ---, each with an explicit conversion rule; the per-class rules, the \texttt{allowDefault} option for pieces without a known weight and the implementation of the \texttt{toGrams} function are detailed in \secref{anexa:grams-detail}.

A single case is worth keeping here, since it is a design decision, not a mere conversion detail: the taste and garnish units (\texttt{TO\_TASTE}, \texttt{TO\_SERVE}, \texttt{GARNISH}) are treated as \texttt{0 g} and excluded from the shopping list. The \texttt{0 g} value is a consciously chosen modeling convention, not an unknown: it honors the intention of the recipe's author (the quantity is deliberately left “to taste”), avoids downstream noise (staple seasonings decremented on every cooking session or listed for shopping are not useful information in an application focused on the waste of bulk food) and keeps a neutral position with respect to dietary intake, which the application is not in a position to prescribe.

The \texttt{0 g} convention does, however, have a known limitation that follows directly from it: staple ingredients seasoned “to taste” (salt, oil, pepper) are not tracked at the quantity level, so their gradual depletion --- through small, repeated uses --- is not signaled to the user. The right solution to this gap is not assigning an arbitrary non-zero value (which would reintroduce the noise described above and would produce, in the absence of a warning infrastructure, data that nothing consumes), but a dedicated mechanism for tracking consumables, based on a state (\textit{in stock} / \textit{running low} / \textit{out of stock}) marked by the user, not on estimated mass. Such a module --- foreshadowed by the manual “running low” marker already present in the pantry interface --- is described as a direction for future development in \secref{sec:directii-viitoare}.

The structural decision this pipeline represents goes beyond the implementation detail: it is a \textit{unification} between two processes that previously used similar but separate heuristics. When generating the initial data sets (\textit{seeds}), an offline cleaning pipeline imported the \textit{TheMealDB} catalog and estimated the missing weights using a manual \texttt{INGREDIENT\_PIECE\_GRAMS} table built for this stage. At runtime, pantry decrementing (\secref{subsec:sesiune-gatit}) and shopping list derivation (\secref{subsec:lista-cumparaturi}) required the same conversions, but through separately written code. By extracting the \texttt{toGrams}/\texttt{fromGrams} functions into a common library (\texttt{libs/shared/utility}) and promoting the manual table to the \texttt{IngredientCatalog.pieceGrams} column, both processes now consume the same source of truth. From an auditing standpoint, a critique of the estimate (for example, “the density of water is a poor approximation for oils”) applies to both consumers at once; from an evolution standpoint, an improvement (a layer of per-category ingredient densities) propagates the same way. This design lesson is formulated as a generalizable architectural pattern in \secref{subsec:pan-patterns}.

The \texttt{toGrams}/\texttt{fromGrams} pair is a textbook application of information hiding \cite{parnas1972}: the changeable design decision --- the concrete way of converting between heterogeneous units (weight factors, approximating volumes through the density of water, the \texttt{pieceGrams} values for piece-type units) --- is isolated behind a stable interface, and the two consumers (the offline \textit{seed} pipeline and the runtime logic) depend exclusively on this interface, not on the implementation. The rule formulated above --- one defensible heuristic, two or more consumers --- restates the criterion according to which a module's interface should reveal as little as possible about its internal mechanism. The benefit is the one described by Parnas: a critique of the estimate (the density of water as a poor approximation for oils) or an improvement (a layer of per-category densities) applies to all consumers at once, because there is a single place where the decision is materialized.

\subsubsection{Deriving the shopping list}\label{subsec:lista-cumparaturi}

On the \texttt{/groceries} route, the shopping list appears as a \textit{set difference on ingredient identity} between the ingredients required by the planned slots and the contents of the pantry. The \texttt{groceries.list} procedure expands each slot with a recipe in the user's plan (all slots with a non-null \texttt{recipeId}) into the ingredients of the corresponding recipe, scales the quantities by the ratio \texttt{slot servings / recipe servings}, aggregates by \texttt{ingredientCatalogId} and completely removes any ingredient for which a row already exists in the pantry --- the removal comparison is made on \textit{presence}, not on quantity. Ingredients with a \textit{flavour-only} unit (\texttt{TO\_TASTE}, \texttt{TO\_SERVE}, \texttt{GARNISH}) are excluded from the list (salt, pepper and garnishes do not appear on the shopping list; \texttt{PINCH} is kept, having a non-zero mass (0.5 g)).

Reconciliation in grams does not happen when filtering the list, but when merging with the pantry: at the end, the \texttt{groceries.addBoughtToPantry} mutation transfers the items checked as “bought” directly into the pantry, merging them with any existing rows for the same ingredient. This merge uses the grams pipeline from \secref{subsec:algoritm-grame} --- it normalizes the quantities to grams, sums them and converts back to the storage unit of the existing row (or of the new item, if the row is missing). Piece-type items without a populated \texttt{pieceGrams} consume \texttt{DEFAULT\_PIECE\_GRAMS = 100} g through the \texttt{allowDefault} option, a \textit{graceful degradation} solution preferable to failing the operation because of a single incomplete ingredient.

The “bought” state of each item is persisted in \gls{redis} (\secref{subsec:rute-h3-redis}), allowing the user to continue the list on different \textit{devices} within the key's expiration window.

\subsection{Methodology and evaluation}\label{ch:evaluare}

\subsubsection{Experimental setup}\label{subsec:setare-experimentala}

All measurements were taken in the \gls{vercel} \textit{Preview} environment (the \texttt{fra1} region, Frankfurt), with \gls{cockroachdb} Serverless and \gls{redis} Cloud in the same region; the stack runs \gls{analog} on top of \gls{nitro}, with \acrshort{ssr} on the first request and \gls{hydration} on the client. The data set comes from deterministic test data sets (\textit{fixtures}, \texttt{libs/shared/mock/}) and from the initialization scripts (\textit{seed}, \texttt{libs/\allowbreak prisma-client/\allowbreak seed-data/\allowbreak catalog.json}, \texttt{recipes.json}).

All the measurement procedures are versioned in the project's repository: the benchmark route is defined in \texttt{apps/\allowbreak web/\allowbreak src/\allowbreak server/\allowbreak routes/\allowbreak api/\allowbreak internal/\allowbreak bench-db.get.ts}, the \acrshort{ci} flows are in \texttt{.github/workflows/}, and the configuration of the \acrshort{e2e} suite is in \texttt{apps/web-e2e/playwright.config.ts}. Each dimension is reproduced through a dedicated script in \texttt{scripts/}: persistence latency through \texttt{npm run bench:preview}, test coverage through \texttt{npm run bench:coverage}, and front-end performance --- the \gls{lighthouse} scores and the image payload --- through \texttt{npm run bench:lighthouse} and \texttt{npm run bench:images}, respectively; the scripts that query the \gls{vercel} environment require the corresponding access credentials.

\subsubsection{Indicators and procedures}\label{subsec:indicatori-proceduri}

\paragraph{Latency at the persistence layer.} The internal \texttt{/api/internal/bench-db} route executes, on each call, \(N=100\) successive read queries against \gls{cockroachdb} and \gls{redis}, respectively, reporting the arithmetic mean, the 50th percentile (\acrshort{p50}) and the 95th percentile (\acrshort{p95}). The cold/warm distinction operates at the level of the \textit{serverless} function invocation: the first call pays for the cold start of the instance and of the \textit{connection pool} to the database, while subsequent calls hit the already warmed-up instance. The cold call is reported separately from the median of the warm calls.

\paragraph{Reliability of the end-to-end (\acrshort{e2e}) tests.} The pass rate of the \gls{playwright} suite on Chromium, measured over the last ten runs on the \texttt{test} branch. WebKit also passes completely, and Firefox runs most of the tests, except a few authentication tests (marked \texttt{@no-firefox}), excluded because of a synchronization race at the level of the \textit{fetch} requests in the local development environment. In \acrshort{ci} only Chromium runs systematically; Firefox and WebKit are covered in manual runs before promotions to \texttt{main}.

\paragraph{Front-end performance.} \gls{lighthouse} is run on the home page of the application published at \texttt{kookio.app}, in the desktop and mobile profiles (the latter with network and CPU \textit{throttling}), reporting the median of several runs for the performance score and for the \gls{lcp}, \gls{cls} and \gls{tbt} metrics. Separately, the payload reduction of the image optimization is quantified: for each of the 594 images in the catalog, the original JPEG file is compared with the \gls{avif} variant delivered by the \gls{vercel} optimizer at the maximum configured width (\(w=960\)) --- a conservative bound, since smaller responsive widths save even more.

\paragraph{Unit test coverage.} Line coverage is reported over the entire source code of the libraries (\textit{all-source}: files not covered by any test are counted as \(0\,\%\), not omitted), aggregated from the per-project \gls{vitest} reports. The \texttt{web} application is excluded from this figure, because its test environment times out under coverage instrumentation; its user flows are validated by the \acrshort{e2e} suite.

\subsubsection{Results}\label{subsec:rezultate}

\begin{table}[htbp]
\centering
\caption{Summary of the results across the three measured dimensions.}
\label{tab:rezultate}
\begin{tabular}{lr}
\toprule
\textbf{Indicator} & \textbf{Value} \\
\midrule
\multicolumn{2}{l}{\textit{Latency at the persistence layer (\(N=100\) probes/call, median over warm calls)}} \\
\quad \gls{cockroachdb} (warm) --- mean & 1.79 ms \\
\quad \gls{cockroachdb} (warm) --- \acrshort{p50} & 1.71 ms \\
\quad \gls{cockroachdb} (warm) --- \acrshort{p95} & 2.07 ms \\
\quad \gls{cockroachdb} (\textit{cold start}) --- \acrshort{p50} / \acrshort{p95} & 2.6 / 6.4 ms \\
\quad \gls{redis} (warm) --- mean & 0.59 ms \\
\quad \gls{redis} (warm) --- \acrshort{p50} & 0.56 ms \\
\quad \gls{redis} (warm) --- \acrshort{p95} & 0.77 ms \\
\midrule
\multicolumn{2}{l}{\textit{Front-end performance (\gls{lighthouse}, median)}} \\
\quad Desktop score & 93 \\
\quad Mobile score (with \textit{throttling}) & 58 \\
\quad \acrshort{lcp} desktop / mobile & 1.5 / 8.3 s \\
\quad \acrshort{cls} (desktop / mobile) & 0.055 / 0 \\
\quad Image payload reduction (594 images, \(w=960\)) & \(-60.5\,\%\) \\
\midrule
\multicolumn{2}{l}{\textit{Coverage and testing}} \\
\quad Line coverage (library source code) & \(82.2\,\%\) \\
\quad Unit tests (\gls{vitest}) & 873 / 96 files \\
\quad \acrshort{e2e} tests (\gls{playwright}) & 14 / 7 flows \\
\quad \acrshort{e2e} pass rate on Chromium & \(100\,\%\) \\
\bottomrule
\end{tabular}
\end{table}

The results summarized in \tabref{tab:rezultate} confirm, within the limits explained in \secref{subsec:discutie-limitari}, the viability of the \gls{pan} architecture along the three evaluation axes: millisecond latencies at the persistence layer, solid front-end performance on desktop (score 93, zero \gls{cls}) with an acknowledged gap in the throttled mobile profile, a reduction of approximately \(60\,\%\) in the image payload across the entire catalog, a coverage of \(82.2\,\%\) of the library code and complete reliability of the critical end-to-end flows on the reference browser.

\subsubsection{Discussion and limitations}\label{subsec:discutie-limitari}

All measurements come from the \textit{Preview} environment, in the \texttt{fra1} region, without concurrent traffic --- so they do not capture the stack's behavior under load or in a multi-region setup. The \acrshort{e2e} suite runs only on Chromium in \acrshort{ci}, and the Firefox and WebKit test matrices are covered manually; automated promotion of the two would increase confidence in \textit{cross-browser} compatibility. In addition, both \gls{cockroachdb} and \gls{redis} run on the free tiers (with usage quotas and an \textit{auto-suspension} cycle for \gls{redis}) and are measured without concurrent traffic, which makes the reported latencies representative of a single-user regime, not of production at scale. Generalizing these results requires dedicated load tests and repeated measurements in a stable production environment. Moreover, these latencies are dominated by the free-tier managed providers --- the network time to the \texttt{fra1} region and the free-quota limitations ---, so they characterize the deployment infrastructure, not the \gls{pan} stack's own \textit{overhead}. Isolating this overhead would require a local, containerized benchmark (for example with Docker), without the network latency of the cloud providers --- a measurement direction left for future work.

Three limitations specific to the new measurements are worth pointing out. First, the persistence latencies vary noticeably between measurement sets --- the \acrshort{p50} median for \gls{cockroachdb} oscillated between approximately \(1.4\) and \(2.6\) ms ---, a consequence of the free tiers, the single region and the variable load; the value reported in \tabref{tab:rezultate} is the median of a set of warm calls, not a constant. Second, the \gls{lighthouse} score on the mobile profile reflects the aggressive \textit{throttling} applied by the tool (slow network and a four-times slower CPU), not the performance perceived on a real device --- the desktop profile remains in the good range. Finally, the \(82.2\,\%\) coverage concerns exclusively the library code, the \texttt{web} application being excluded for the reasons explained in \secref{subsec:indicatori-proceduri}.
% 3-implementation.tex ends here
